% ============================================================================
% 20_results.tex  --  Methods & results  (~2 pages, the core of the report)
%
% GOAL: What you computed and what it shows. Include the key figure(s) for your
%       part ONLY (see the Code/figs column for your part in
%       ../../reports/README.md). Regenerate fresh PDFs with the block runner,
%       e.g.  python src/block1.py --plots 6   (venv /opt/python-envs/myenv),
%       and copy them into a local figs/ folder. Do NOT use showcase.py gallery
%       images (shared). Report concrete numbers, not just plots.
% ============================================================================

\section{Methods and Results}
\label{sec:results}

\subsection{Circuit-level implementation and verification}

\ai{All circuit results use a chain of $L=4$ sites, hopping strength $t=1$,
pairing strength $\Delta=1$, and the parity-constrained VQE and edge string from
Part 4.  The chemical potential $\mu$ is varied to cross the phase transition.
The preparation circuit contains single-qubit $R_Y$ rotations and linear CNOT
gates~\cite{peruzzo2014vqe}.  Before simulation it is written in the fixed gate
set \texttt{[rz,sx,cx]} without removing gates.  An error assigned to
\texttt{cx} therefore acts after every CNOT.  The output density matrix $\rho$
gives the edge-string value
$\operatorname{Tr}(O_{\mathrm{edge}}\rho)$.}

\ai{The implementation was checked with a separate calculation that applies
each gate and each depolarizing channel directly, without using Aer's noise
machinery.  Across 41 values of $\mu/t$, the two calculations differ by only
about $10^{-15}$.  Figure~\ref{fig:verification} also compares the deliberately
strong CNOT depolarizing strength $p_{\mathrm{cx}}=0.10$ with thermal relaxation
at $T_1=80\,\mu\mathrm{s}$, $T_2=60\,\mu\mathrm{s}$ and CNOT duration
$\tau_{\mathrm{cx}}=300\,\mathrm{ns}$.  The relaxation probability per gate is
only $1-e^{-\tau_{\mathrm{cx}}/T_1}\simeq0.0037$, so the thermal curve remains
closer to the ideal result.}

\begin{figure}[H]
  \centering
  \includegraphics[width=0.98\linewidth]{block4_week9_verification.pdf}
  \caption{Left: the Aer \texttt{NoiseModel} edge string and the independent
  exact gate-by-gate reference agree to $\sim10^{-15}$ across the $\mu$ sweep.
  Right: depolarizing versus thermal $T_1/T_2$ noise at matched gate placement.}
  \label{fig:verification}
\end{figure}

\subsection{The optimal depth is an expressibility threshold}

\ai{At the topological sweet spot $\mu=0$, the noiseless edge string is only
$\simeq0.004$ for one repetition ($r=1$), but reaches approximately one at
$r=2$ and changes little in deeper circuits.  The first circuit is simply too
shallow to represent the state.  With $p_{\mathrm{cx}}=0.05$, the measured
signal peaks at $r^*=2$; every further repetition adds three noisy CNOTs without
improving the ideal state.  Figure~\ref{fig:depth} shows these two effects
separately.}

\begin{figure}[H]
  \centering
  \includegraphics[width=0.72\linewidth]{block4_week9_depth_optimum.pdf}
  \caption{For $L=4$, the noiseless edge string has an expressibility threshold:
  $r=1$ fails, $r=2$ succeeds, and deeper circuits mainly accumulate CNOT noise.}
  \label{fig:depth}
\end{figure}

\ai{Angle errors provide a second check.  A random shift $\delta\theta$ rotates
the ideal state away from the optimum while its purity stays at one.  CNOT
depolarization reduces the same edge string while the purity falls.  The first
error can in principle be corrected by recalibrating the angles; the second
cannot be removed by changing one set of angles.}

\subsection{Noisy optimization and backend calibration}

\ai{This experiment fixes $r=3$, which does not replace the $r=2$ optimum found
above.  The earlier test asks which depth gives the largest measured signal for
$p_{\mathrm{cx}}=0.05$.  Here the purpose is different: the depth is held fixed
one layer above the representation threshold so that only the response of the
optimizer to two noise channels is compared.  Thus $r=3$ is a controlled test
setting, not a recommended operating depth; its extra CNOTs also explain why its
absolute noisy values are lower than an $r=2$ run would be.}

\ai{The noisy-cost calculation uses $L=4$, $r=3$, $\mu=0$, parity-penalty
strength $\lambda=0.1$, and the COBYLA optimizer.  It starts from the noiseless
angles and also tries several random starting points.  A channel is
\emph{unital} if it leaves the fully mixed state unchanged; it does not pull the
system toward one preferred state.  Depolarizing noise is unital, and optimizing
again changes the result very little.  A \emph{non-unital} channel does prefer a
state.  Amplitude damping is non-unital because it drives qubits toward
$|0\rangle$.  Under this channel, re-optimization can increase the
\emph{measured} edge string while moving the corresponding noiseless circuit
away from the exact ground state.  Table~\ref{tab:noisy} gives selected values
from the seed-7 sweep.}

\begin{table}[H]
  \centering
  \small
  \begin{tabular}{llcccc}
    \toprule
    channel & strength & \multicolumn{2}{c}{$|\langle O_{\mathrm{edge}}\rangle|$ noisy}
      & \multicolumn{2}{c}{fidelity to exact GS}\\
    \cmidrule(lr){3-4}\cmidrule(lr){5-6}
      & & fixed $\theta^*$ & noise-aware & fixed $\theta^*$ & noise-aware\\
    \midrule
    depolarizing & $0.08$ & $0.502$ & $0.501$ & $0.998$ & $1.000$\\
    amplitude damping & $0.08$ & $0.422$ & $0.565$ & $0.998$ & $0.994$\\
    amplitude damping & $0.12$ & $0.266$ & $0.337$ & $0.998$ & $0.966$\\
    \bottomrule
  \end{tabular}
  \caption{Best-case versus noise-aware optimization at $L=4$, $r=3$, $\mu=0$.
  Under amplitude damping the measured edge string rises while the noiseless
  fidelity to the exact ground state falls at strong damping.}
  \label{tab:noisy}
\end{table}

\ai{Finally, \texttt{NoiseModel.from\_backend} replaces the uniform test channel
with a recorded \texttt{FakeManilaV2} calibration.  Its CNOT error rates vary
from about $0.009$ to $0.014$.  At $\mu=0$ the ideal, backend-calibrated, and
uniform $p_{\mathrm{cx}}=0.05$ edge strings are $0.99$, $0.87$, and $0.67$.
The uniform model is therefore a deliberately harsh comparison, not a device
prediction.  The backend result is still optimistic: its readout errors are not
used because the simulation evaluates
$\operatorname{Tr}(O_{\mathrm{edge}}\rho)$ without measurement instructions.
The data also come from a recorded calibration, not a live hardware run.}

\ai{Together these tests show that realistic noise primarily reduces the
contrast of the non-local diagnostic.  Circuit depth should be chosen by
noiseless expressibility validation followed by noisy evaluation, and any
noise-aware optimization must report a state-quality check rather than only the
observable it was able to improve.  These conclusions concern the $L=4$ NISQ
circuit; their large-$L$ consequences are treated separately in Part 6.}
